% ============================================================================
%  SECCIÓN 4.10: PRUEBAS FUNCIONALES Y DIFICULTADES DEL AVANCE
% ============================================================================

\section{Pruebas funcionales y dificultades del avance}

\subsection{Pruebas manuales de la l\'ogica implementada}

Siguiendo la gu\'ia del avance, se realizaron pruebas manuales sobre los
formularios, las validaciones y la funcionalidad principal interactiva.
Las pruebas se ejecutaron en un dispositivo f\'isico Android mediante
Expo Go y cubren los siete casos propuestos por la gu\'ia:

\begin{table}[H]
  \centering
  \caption{Resultados de las pruebas funcionales del avance}
  \label{tab:pruebas-logica}
  \contenidoConFuente{%
    \begin{tablaAPA}
      \begin{tabular}{@{}p{0.8cm}p{3.6cm}p{2.0cm}p{8.6cm}@{}}
        \toprule
        \textbf{ID} & \textbf{Prueba} & \textbf{Resultado} &
        \textbf{Observaciones} \\
        \midrule
        P1 & Formulario vacío & Aprobado & Al presionar el botón sin
        completar campos, cada campo muestra su mensaje de obligatoriedad
        debajo en color rojo; el envío no se ejecuta. \\
        \addlinespace
        P2 & Formulario parcialmente lleno & Aprobado & Solo los campos
        incompletos o inválidos son señalados; los válidos conservan el
        borde normal. \\
        \addlinespace
        P3 & Datos incorrectos & Aprobado & Con un email sin formato
        válido o una contraseña menor a 6 caracteres, Zod muestra los
        mensajes específicos de cada regla. \\
        \addlinespace
        P4 & Datos correctos & Aprobado & El botón se deshabilita y
        muestra un indicador de carga durante el envío; al finalizar se
        informa que la cuenta fue creada. \\
        \addlinespace
        P5 & Acción principal (favorito) & Aprobado & El toggle añade o
        elimina el escenario en Supabase y la lista de favoritos se
        actualiza tras invalidar el caché de React Query. \\
        \addlinespace
        P6 & Cambio de estado & Aprobado & El botón cambia de estado
        vacío a lleno; mientras dura la operación queda deshabilitado,
        evitando toques duplicados. \\
        \addlinespace
        P7 & Caso de error & Aprobado & Con credenciales inválidas se
        muestra la caja de error del servidor; con un escenario
        inexistente se muestra el estado vacío en lugar de una pantalla
        en blanco. \\
        \bottomrule
      \end{tabular}
    \end{tablaAPA}%
  }{Elaboración propia en base a las pruebas realizadas en dispositivo
  físico.}
\end{table}

\subsection{Dificultades encontradas en el avance}

Durante el desarrollo de este avance se encontraron las siguientes
dificultades, junto con las soluciones aplicadas:

\begin{table}[H]
  \centering
  \caption{Dificultades del avance y soluciones aplicadas}
  \label{tab:dificultades-avance4}
  \contenidoConFuente{%
    \begin{tablaAPA}
      \begin{tabular}{@{}p{5.5cm}p{9.5cm}@{}}
        \toprule
        \textbf{Problema} & \textbf{Solución} \\
        \midrule
        El registro con correos ficticios no permitía iniciar sesión:
        Supabase creaba la cuenta pero exigía confirmación por correo,
        imposible de completar con un buzón inexistente. & Se identificó
        que el proyecto alojado tiene activada la opción \textit{Confirm
        email}. Se implementó la verificación mediante código de un solo
        uso (\textit{OTP}) enviado al correo real, con pantalla dedicada
        para ingresarlo y opción de reenvío. \\
        \addlinespace
        El servidor de correo integrado de Supabase limita el envío a
        aproximadamente dos mensajes por hora, insuficiente para las
        pruebas del equipo. & Se ajustaron los límites de autenticación
        disponibles en el panel y se documentó la configuración de un
        servidor SMTP externo gratuito para despliegues con mayor
        volumen. \\
        \addlinespace
        La lista de escenarios no reflejaba los cambios al agregar o
        quitar un favorito hasta reiniciar la aplicación. & Tras cada
        mutación se invalidan las consultas de React Query asociadas
        (\textit{favorites} y \textit{favorite}), forzando la
        resincronización de la interfaz. \\
        \addlinespace
        La búsqueda filtraba mientras el usuario aún escribía,
        generando filtrados intermedios innecesarios. & Se introdujo un
        retardo (\textit{debounce}) de 300 ms con limpieza del temporizador
        al desmontar el componente, evitando fugas de memoria. \\
        \addlinespace
        Los botones permitían toques repetidos durante operaciones
        asíncronas, duplicando registros de favoritos. & Se incorporaron
        estados de envío (\textit{isSubmitting}, \textit{isToggling}) que
        deshabilitan el botón y muestran un indicador de carga. \\
        \addlinespace
        Al consultar un escenario inexistente la pantalla quedaba en
        blanco sin retroalimentación. & La función de consulta retorna
        \textit{null} ante el código \textit{PGRST116} y la interfaz
        muestra el componente de estado vacío con mensaje informativo. \\
        \bottomrule
      \end{tabular}
    \end{tablaAPA}%
  }{Elaboración propia en base a las dificultades encontradas durante el
  desarrollo.}
\end{table}
